\documentclass[onecolumn, draftclsnofoot,10pt, compsoc]{IEEEtran}
\usepackage{graphicx}
\usepackage{url}
\usepackage{setspace}
\usepackage[parfill]{parskip}
\usepackage{geometry}
\geometry{textheight=9.5in, textwidth=7in}
\usepackage[pdf]{pstricks}
\usepackage{pst-gantt}

% 1. Fill in these details
\def \CapstoneTeamName{Aggregators}
\def \CapstoneTeamNumber{71}
\def \GroupMemberOne{Carter Olsen}
\def \GroupMemberTwo{Megan Liles}
\def \GroupMemberThree{Aalok Borkar}
\def \GroupMemberFour{Race Stewart}
\def \CapstoneProjectName{News Aggregator Web Development}
% \def \CapstoneSponsorCompany{	Cheap Robots, Inc}
\def \CapstoneSponsorPerson{Joseph Louis}

% 2. Uncomment the appropriate line below so that the document type works
\def \DocType{		%Problem Statement
				Group Problem Statement
				%Technology Review
				%Design Document
				%Progress Report
				}
			
\newcommand{\NameSigPair}[1]{\par
\makebox[2.75in][r]{#1} \hfill 	\makebox[3.25in]{\makebox[2.25in]{\hrulefill} \hfill		\makebox[.75in]{\hrulefill}}
\par\vspace{-12pt} \textit{\tiny\noindent
\makebox[2.75in]{} \hfill		\makebox[3.25in]{\makebox[2.25in][r]{Signature} \hfill	\makebox[.75in][r]{Date}}}}
% 3. If the document is not to be signed, uncomment the RENEWcommand below
% \renewcommand{\NameSigPair}[1]{#1}

%%%%%%%%%%%%%%%%%%%%%%%%%%%%%%%%%%%%%%%
\begin{document}
\begin{titlepage}
    \pagenumbering{gobble}
    \begin{singlespace}
    	\includegraphics[height=4cm]{ProblemStatment/coe_v_spot1.eps}
        \hfill 
        % 4. If you have a logo, use this includegraphics command to put it on the coversheet.
        %\includegraphics[height=4cm]{CompanyLogo}   
        \par\vspace{.2in}
        \centering
        \scshape{
            \huge CS Capstone \DocType \par
            {\large\today}\par
            \vspace{.5in}
            \textbf{\Huge\CapstoneProjectName}\par
            \vfill
            {\large Prepared for}\par
            % \Huge \CapstoneSponsorCompany\par
            \vspace{5pt}
            {\Large\CapstoneSponsorPerson\par}
            {\large Prepared by }\par
            Group \CapstoneTeamNumber\par
            % 5. comment out the line below this one if you do not wish to name your team
            \CapstoneTeamName\par 
            \vspace{5pt}
            {\Large
                \GroupMemberOne\par
                \GroupMemberTwo\par
                \GroupMemberThree\par
                \GroupMemberFour\par
            }
            \vspace{20pt}
        }
        \begin{abstract}
        % 6. Fill in your abstract    
        	What is one of the biggest problems in our social media today?
        	Bias.
        	It is impossible to read about current events without bias clouding  the factual evidence. 
        	Readers are often mislead about the severity of events and are told factual errors about current events.
        	Today's media has high levels of inaccuracies, an overemphasis on emotion, poor coverage of all important events, and a bigger focus on profits rather than educating the public.
        	One solution is to create a website that gathers news articles for the current events from multiple sources on the internet.
        	Currently, users only get to read from one point of view on a website and if they want multiple perspectives they need to visit multiple websites.
        	This solution will use machine-learning to gather articles from multiple perspectives and aggregate them into one place.
        \end{abstract}     
    \end{singlespace}
\end{titlepage}
\newpage
\pagenumbering{arabic}
\tableofcontents
% 7. uncomment this (if applicable). Consider adding a page break.
%\listoffigures
%\listoftables
\clearpage

% 8. now you write!
\section{Description}

There are several problems that exist currently in today’s media coverage.
 The biggest problem in today’s media is bias.
 News sources exhibit their beliefs, morals and values throughout their media and try to cast it upon their audience.
 This is a problem because readers and listeners are not getting a factual description or story of current events they are absorbing.
 Not only do news companies display bias in their articles and other media but they provide inaccurate factual evidence, an increase in sentiment causing stories to be more emotionally heightened than usual, and a strong motivation for profit rather than educating the audience.


A majority of Americans believe that their news source does not get all of their facts straight and many audiences find factual errors in their news on a daily basis.
 Oftentimes news sources do not get to witness the event they are reporting on and end up providing inaccurate facts or evidence as they have only sourced it from another news source.
 

Sensationalism is another common problem in today’s media.
 Many news companies relish in reporting emotionally-hard events such as kidnapping, rape, murder, robberies, natural  disasters, accidents, scandals and more.
 Many audiences favorite these types of stories.
 Unfortunately, news reporters will oftentimes play into these emotional stories and exaggerate the severity of them to gain popularity.
 They also overlook the positive news stories and will report a bad story over a good one.
 This misleads the general public about what things are happenings around them and can convince the audience that there is always more negative events rather than positive ones.


Remarkably, rarely do news companies ever correct their own errors in the media.
 Many organizations often are reluctant to do so in fear of discrediting their own employees.
 Most Americans believe that the media should have guidelines to report their errors.
 This again is a problem as it misleads users about what is really happening in the world around them.
 People need to know the facts as it can influence their thinking in various ways.
 The media influences all sorts of decisions whether it is about voting, preparing for natural disasters, buying a new  home, looking for a  new job, purchasing stocks, or even trying a new restaurant.


Another problem that exists in media is the selection of coverage on important events.
 Oftentimes, news organizations will pick and choose between which stories to cover.
 One important example is the fire in the Amazon.
 There was very little media coverage about this current event and many other environmental issues.
 News organizations will purposely not report on new bills being passed that could be very controversial.
 Foreign countries, education, shootings, wars and many other events are also left out of media coverage depending on which news source users visit.

\section{Proposed Solution}

The general public is in need of a news source that is neutral and free of bias.
 Audiences need a news source that aggregates all of the current events from multiple perspectives and hosts them in one location so they can read from various views and get a more accurate description of the event.
 The proposed solution will consists of a website that aggregates news articles using machine learning methods to gather these stories from RSS feeds by topic and host them all in one central location on a website.
 These stories will be gathered from specified news sources and displayed on one website so users will be able to read about a current event from multiple perspectives.
 This website will also have a feature that allows users to be able to share these posts on popular social media platforms.
 It is not easy to control what news organizations write about and how they overemphasize the emotion in media, make factual errors or overlook important topics.
 However, these problems in the media can be brought to attention when users start to read from the aggregated news website and discover the discrepancies between stories.
 This will allow them to at the very least make an informed and educated opinion about current events and bring to light the uncertainties of the inaccurate details throughout media.

\section{Performance Metrics}
There are several components that measure the successful completion of this project.
 The first being usability.
 Are users able to successfully navigate the website without any flaws or errors? Is the website intuitive to navigate and use? How do users feel when using the website? Is there any distractions on the website that prohibit the user from experiencing the website in the intended way? Is the website organized in a way that makes sense? Inclusive design is another important factor to measure when doing web development.
 The user needs to be able to remember how to use the website easily and not have to teach themselves each time they visit the site.
 The site must be easy to use and also must provide ways to allow underrepresented populations use the site (blind users, deaf users, users with low motor skills, etc.
) Most importantly the website must do the intended job of aggregating news from multiple sources.
 Are multiple sources equally represented or is one news source overpowering others? Does the website allow users to search a current event and is it efficient? Lastly, the website needs to keep content up to date and provide users with new events, not just top stories.
 



\end{document}